% !PRINTED 53
\weberpage{53}
%<<CONT

\gesperrt{Die $\mu$te Derivirte von der $\nu$ten Derivirten ist die $(\nu+\mu)$te Derivirte der ursprünglichen Function.}

Man erhält also die sämmtlichen höheren Derivirten, indem man nach der Regel (3) aus jeder vorangehenden die erste Derivirte bildet:
\begin{align}
 f(x)&=a_0x^n+a_1x^{n-1}+a_2x^{n-2}+a_3x^{n-3}+\cdots,\notag\\
 f'(x)&=na_0x^{n-1}+(n-1)a_1x^{n-2}+(n-2)a_2x^{n-3}+\cdots, \tag{9}\\
 f''(x)&=n(n-1)a_0x^{n-2}+(n-1)(n-2)a_1x^{n-3}+\cdots. \notag
\end{align}

Eine etwas einfachere Form nehmen diese Derivirten an, wenn man sich einer anderen Bezeichnungsweise bedient, die häufig im Gebrauch und für gewisse Zwecke fast unentbehrlich ist, die wir im Anschluss hieran besprechen wollen.

Es liegt wegen der Unbestimmtheit der Coëfficienten $a_0,a_1,\ldots,a_n$ offenbar keine Beschränkung darin, wenn wir eine ganze rationale Function $n$ten Grades so darstellen:
\begin{equation}
 f(x)=a_0x^n+B_1^{(n)}a_1x^{n-1}+B_2^{(n)}a_2x^{n-2}+\cdots+a_n,
\tag{10}
\end{equation}
oder ausführlich:
\begin{equation}
 f(x)=a_0x^n+na_1x^{n-1}+{n(n-1)\over1\cdot2}a_2x^{n-2}+\cdots.
\tag{11}
\end{equation}
Wenn eine Function $f(x)$ so dargestellt ist, werden wir sagen, sie sei ``\gesperrt{mit den Binomialcoefficienten geschrieben}''.

Grössere Uebereinstimmung zeigen hierdurch bereits die Formeln (9), die dann so lauten:
\begin{align}
 f(x)&=a_0x^n+na_1x^{n-1}+{n(n-1)\over1\cdot2}a_2x^{n-2}+\cdots,\notag\\
 {1\over n}f'(x)&=a_0x^{n-1}+(n-1)a_1x^{n-2}+{(n-1)(n-2)\over1\cdot2}a_2x^{n-3}+\cdots, \tag{12}\\
 {1\over n(n-1)}f''(x)&=a_0x^{n-2}+(n-2)a_1x^{n-3}+{(n-2)(n-3)\over1\cdot2}a_2x^{n-4}+\cdots,\notag
\end{align}
worin die rechten Seiten alle auch mit den Binomialcoefficienten geschrieben erscheinen. Wir werden später den Nutzen dieser Bezeichnungsweise noch weiter kennen lernen, müssen aber schon hier hervorheben, dass die Wahl der einen oder anderen Dar-
%CONT>>
